\begin{figure*}[t]
\centering
\begin{tikzpicture}[
    node distance=1.8cm and 1.8cm,
    >=Latex,
    box/.style={
        draw,
        rounded corners,
        align=center,
        minimum width=4.1cm,
        minimum height=1.5cm,
        font=\small
    },
    smallbox/.style={
        draw,
        rounded corners,
        align=center,
        minimum width=3.6cm,
        minimum height=1.1cm,
        font=\small
    },
    arrow/.style={->, thick}
]

\node[box] (cal) {
    \textbf{Calibration acquisition}\\
    $\sim 2\times 10^4$ raw events
};

\node[smallbox, below=0.45cm of cal] (caluse) {
    Photopeak / ToT selection\\
    Baseline-quality cuts\\
    LED threshold choice\\
    Fixed timing-offset estimate
};

\node[box, right=3.2cm of cal] (train) {
    \textbf{Training acquisition}\\
    $\sim 10^5$ raw events
};

\node[smallbox, below=0.45cm of train] (trainuse) {
    Fixed preprocessing\\
    3-fold cross-validation\\
    Model / hyperparameter selection by CTR\\
    Final fit on full training set
};

\node[box, right=3.2cm of train] (test) {
    \textbf{Test acquisition}\\
    $\sim 10^5$ raw events
};

\node[smallbox, below=0.45cm of test] (testuse) {
    Fixed preprocessing\\
    Final model evaluation\\
    CTR and RMSE\\
    Bootstrap uncertainty
};

\draw[arrow] (cal.east) -- node[above, font=\small, align=center] {fixed selection rules\\and LED calibration} (train.west);
\draw[arrow] (train.east) -- node[above, font=\small, align=center] {selected model\\and fitted parameters} (test.west);

\draw[arrow] (cal) -- (caluse);
\draw[arrow] (train) -- (trainuse);
\draw[arrow] (test) -- (testuse);

\end{tikzpicture}
\caption{Schematic overview of the experimental protocol. The calibration acquisition is used to determine fixed event-selection and LED-calibration criteria. Model selection and fitting are performed on the independent training acquisition. Final performance is assessed on the independent test acquisition.}
\label{fig:workflow}
\end{figure*}